\documentclass[25pt,landscape,headrule,footrule,dvips]{foils}

\leftheader{\makebox[10mm]{\raisebox{0mm}[0mm][0mm]{
\includegraphics[width=50mm]{ntulogo}}}}
\MyLogo{\shortstack{PACLIC-25}}
\newcommand{\section}[1]{\myslide{}{\begin{center}\Huge \emp{#1}\end{center}}}
\newcommand{\Bad}{\emp{\raisebox{0.15ex}{\ensuremath{\mathbf{\otimes}}}}}

\usepackage{CJKutf8}
\usepackage{CJK}
\begin{CJK}{UTF8}{min} 
\title{\vspace{-3ex}\emp{Nanyang Technological University
     \\	Multilingual Corpus \\ (NTU-MC)}}
\author{ \bl{Francis Bond}, \bl{Liling Tan}, \bl{Shan Wang}, \ldots
  \\ School of Humanities and Social Sciences
  \\ \textbf{Nanyang University of Technology}
  \\ \url{bond@ieee.org}}
\end{CJK}
\date{2014-03-05}

\usepackage{graphicx}
\usepackage{mathpazo}
% \usepackage[T1]{fontenc} 
% \usepackage[utf8]{inputenc}
%\usepackage[vietnam, USenglish]{babel}

%\usepackage[dvips]{color}
\usepackage[dvips]{xcolor}
\newcommand{\ul}{\underline}
\newcommand{\into}{\ensuremath{\rightarrow}\xspace}
\newcommand{\tot}{\ensuremath{\leftrightarrow}\xspace}

\newcommand{\emp}[1]{\textcolor{red}{#1}}
\newcommand{\fd}[1]{\texttt{#1}}
\newcommand{\com}[2][red]{\hfill(\textcolor{#1}{#2})}
\newcommand{\bl}[1]{\textcolor{blue}{#1}}
\usepackage{url}
\newcommand{\wn}{\url}
\usepackage[e,f,j]{mtg2e}

\newcommand{\myslide}[1]{\foilhead[-25mm]{\raisebox{12mm}[0mm]{\emp{#1}}}}

\usepackage{pifont}
\renewcommand{\labelitemi}{\textcolor{violet}{\ding{227}}}
\renewcommand{\labelitemii}{\textcolor{purple}{\ding{226}}}

%\usepackage[dvips]{hyperref}
\begin{document}
\begin{CJK}{UTF8}{min}
\newcommand{\zh}[1]{\CJKfamily{gbsn}#1\CJKfamily{min}}
\newcommand{\vi}[1]{\begin{otherlanguage}{vietnam}#1\end{otherlanguage}}
\maketitle

%%%%%%%%%%%%%%%%%%%%%%%%%%%%%%%%%%%%%%%%%%%%%%%%%%%%%%%%%%%%%%%%%%%%%%%%%%%
\myslide{Overview}
\begin{itemize}
\item \textbf{Motivation for NTU-MC}
\item \textbf{Corpus Construction}
\item \textbf{Annotation and Evaluations}
\item \textbf{Future work and Conclusions}
%\item \textbf{Conclusion}
\end{itemize}
%%%%%%%%%%%%%%%%%%%%%%%%%%%%%%%%%%%%%%%%%%%%%%%%%%%%%%%%%%%%%%%%%%%%%%%%%%%

\myslide{Motivation for NTU-MC}
\begin{itemize}\addtolength{\itemsep}{-0.8ex}
\item Parallel Corpora have proved very useful for various research:
  machine translation, lexicography, \ldots
  \begin{itemize}
\item Many European parallel Corpora (e.g. Europarl, MULTEXT-East).
\item Various  Asian Language-English bitexts 
\\ EMILLE corpus, LCMC (Lancaster)
\item A handful of parallel Asian Language corpora (not open)
  \\ CEJ Patent corpus, NICT Multilingual corpus (CEJ)
  \\ King MT-001 CEKJ corpus, BTEC 
\end{itemize}
\item NTU-MC is building  \textbf{an open corpus of linguistically diverse languages}.
\end{itemize}

%%%%%%%%%%%%%%%%%%%%%%%%%%%%%%%%%%%%%%%%%%%%%%%%%%%%%%%%%%%%%%%%%%%%%%%%%%%

\myslide{Motivation for NTU-MC}
\begin{itemize}\addtolength{\itemsep}{-0.8ex}
\item \textbf{Why a linguistically diversed corpus? \\ --- Computationally}
  \begin{itemize}
  \item Computationally, language diversity allows us to disambiguate
    one language with the other
  \item NTU-MC provides a source of linguistically diverse texts to
    test how well techniques work on different languages
  \item Gives us a test-bed for various theories of representation
      \end{itemize}
\end{itemize}

\myslide{Motivation for NTU-MC}
\begin{itemize}\addtolength{\itemsep}{-0.8ex}
\item \textbf{Why a linguistically diversed corpus? \\--- Linguistically}
	\begin{itemize}
	\item NTU-MC creates a
          platform for traditional linguistic researches on language
          uses in Singapore(e.g. Translation studies, corpus
          linguistics, etc.)
        \item We can compare typological differences on equivalent text
        \item We can perform quantitative contrastive linguistics
        \item We can link to grammars/wordnets for dynamic corpus creation
	\end{itemize}
\end{itemize}
%%%%%%%%%%%%%%%%%%%%%%%%%%%%%%%%%%%%%%%%%%%%%%%%%%%%%%%%%%%%%%%%%%%%%%%%%%%

\myslide{Corpus Construction}
\begin{itemize}\addtolength{\itemsep}{-0.8ex}
\item \textbf{Approach:} The NTU-MC adopts an opportunistic data collection approach.
  \begin{itemize}
  \item Singapore Tourism Websites from \url{www.yoursingapore.com}
    \\ Thanks to the tourist board for permission to redistribute
  \item Singapore Medical Tourism Pages
  \item Sherlock Holmes short stories
  \item Open Source Essays (the Cathedral and the Bazaar)
  \item News (Kyoto Corpus translated by NICT).
  \end{itemize}
\end{itemize}

\myslide{NTU Multilingual Corpus}

{\footnotesize
\hspace{-15mm}
\begin{tabular}{llrrrrrr}\addtolength{\columnsep}{-1ex}
\textbf{Genre} &
\textbf{Text} &
\multicolumn{4}{c}{\textbf{Sentences}} 
& \textbf{Words} & \textbf{Concepts} \\
&& \textbf{Eng} &
\textbf{Cmn} &
\textbf{Jpn} &
\textbf{Ind} &
\textbf{Eng} &
\textbf{Eng} \\
\hline
Story &
Dancing Men &
599 &
606 &
698 &
$-$ &
11,200 & 
5,300 \\
~
 &
Speckled Band &
599 &
612 &
702 & 
$-$ &
10,600 &
4,700
\\
% (- 21818 11198) 10620
Essay &
Cathedral and the Bazaar &
769 &
750 &
773 &
$-$ & 
18,700 & 
8,800\\
News &
Mainichi News&
2,138 &
2,138 &
2,138 &
$-$ &
55,000 &
23,200\\
Tourism &
Your Singapore (web site)&
2,988 &
2,332 &
2,723 &
2,197 &
74,300 &
32,600  
\end{tabular}
}

\begin{itemize}
\item All redistributable (except Mainichi: the WSJ of Japan)
\item All fun to read (except Mainichi)
\item Many translations exist (mainly public domain)
\item Different genres
\end{itemize}

\myslide{Languages and Status}

\begin{small}
  \centering
  \begin{tabular}{lrllllllll}
    Genre & Sents        & \multicolumn{5}{c}{Languages} & Annotation & Comment \\
          & (k) &   cmn & eng & jpn & ind & other\\
\hline
    Tourism1 & 3.3 & C &  \textbf{C} & P  & .65C & kor, vie, tha 
    &   & \\
    Story    & 1.2  & C &  \textbf{C} & C & --- &  
    & xling  & Other corpora \\
    Essay    & 0.8 & C & \textbf{C} & C & --- & many 
    & HPSG (eng) \\
    News     & 2.0  & C & C & \textbf{.5C}  & --- &
    & Sense, Dep (jpn) & $10^6$ words\\
    \hline
    Tourism2 & 3.8 &  P & \textbf{P} & P & P & arb, vie\\
    Semcor & 14.2 &   --- & \textbf{C} & .4C & ---& & HPSG, Penn\\ % (/ 58.0 151)0.3841059602649007
    Gloss  & 112.0 & --- & \textbf{C} & P & ? & alb, spa, \ldots\\
  \end{tabular}
  % \caption{Current State of the Tagged Texts from the NTU-MC}
  % \label{tab:ntu-mc-now}
  \begin{flushleft} \small
    \textbf{Bold} is source language
    \\ $^*$ less for non-English $\approx 2$k
    \\ P: POS 
    \\ C: Concepts (WN synsets), implies POS 
  \end{flushleft}
\end{small}




\myslide{Tourism Corpus}
\begin{itemize}\addtolength{\itemsep}{-0.8ex}
\item \textbf{Size:} 375,00 words (15,000 sentences)
\item \textbf{Language Diversity:} 6 languages from 6 language families
  \begin{itemize}
  \item English (en): Indo-European
  \item Indonesian (id): Austronesian
  \item Japanese (ja): Japonic
  \item Korean (ko): Language Isolate
  \item Vietnamese (vi): Autro-Asiatic 
  \item Chinese (zh): Sino-Tibetan
  \end{itemize}
\item Lots of Singaporean words --- nice for locals
\end{itemize}

\myslide{Stories}
\begin{itemize}
\item Widely studied
\item Many translations
\item Offers to tag in other languages (Thai, Russian, German, Spanish, \ldots)
\item Other annotation
  \begin{itemize}
  \item Information Structure (Sanghoun)
  \item Negative Scope (*SEM)
  \item Rhetorical Argument (Gregoire)
  \end{itemize}
\end{itemize}

\myslide{Essay}
\begin{itemize}
\item Many translations
\item Support of the author
\item Ideologically Sound (promotes open source)
\item Very convoluted syntax --- very hard to parse
\item Gets visibility outside linguistics
\end{itemize}

\myslide{News}
\begin{itemize}
\item Some word alignment
  \\ annoyingly not in the set we chose --- NICT lost the full data
\item Maybe we should give this up?
\item But has lots of useful annotation in Japanese
  \begin{itemize}
  \item Dependencies
  \item Argument Structure
  \item GT tags
  \end{itemize}
\end{itemize}

\myslide{Tagging Goals}
\begin{itemize}
\item Corrected segmentation/lemmatization
\item Wordnet Senses for Essay, Stories, Tourism, News x (cmn, eng, jpn, ind)
\item Pronouns (added to wordnet)
\item Cross-lingual sense links
\item MRS (tomorrow)
\end{itemize}

\myslide{Accessibility Goals}
\begin{itemize}
\item Download (as KAF, SQL)
\item Online lookup
  \begin{itemize}
  \item Search by word/concept/pos (in multiple languages)
  \item Display matching text/word/concept
  \end{itemize}
\item Package for adding a new corpus/tagging
\item Text2Concepts (Kaffinator) for cmn, jpn, eng, ind
\item Statistics of interest package
\end{itemize}

\myslide{Interesting Things}
\begin{itemize}
\item Some POS extensions
  \begin{itemize}
  \item auxiliary verbs
  \item reduplication
  \item -nya, ter-, di-, -lah
  \item sahen v/n
  \end{itemize}
\end{itemize}

\myslide{Data Format}

\begin{itemize}
\item Each genre+language is a separate file
\item xlinks are also separate
\item Unique sentence ID across all files
  \begin{center}\small
    \begin{tabular}{rlll}
      Numbers & Genre & Corpus & Comment \\
      \hline
      1      & essay & catb &  \\
      10,000 & story & danc &  \\
      11,000 & story & spec &  \\
      60,000 & news & kc & and links to original numbers \\
      100,000 & tourism & yoursing &  \\
      110,000 & tourism & med & not annotated \\
      120,000 & gloss & wngloss & not annotated  \\
      300,000 & semcor & semcor & not annotated  
    \end{tabular}
  \end{center}

\end{itemize}


\myslide{Monolingual Corpus}
\begin{itemize}
\item META
\item SENT
\item SUBCORP
\item WORD
\item CONCEPT
\end{itemize}
\myslide{Meta}

\begin{table}[htpb]
  \centering
  \begin{tabular}{llll}
  \textbf{row}   &  \textbf{type} & \textbf{data} & \textbf{comment} \\
  \hline
  corpus & string & corpus name & name without \texttt{.db}\\
  title & string & display name & \\
  language & string & text language & iso 3 letter  \\
\hline
\end{tabular}
\caption{\texttt{meta}: Meta for words}
\end{table}

\myslide{Subcorpus Info}
  \begin{tabular}{llll}
  \textbf{row}   &  \textbf{type} & \textbf{data} & \textbf{comment} \\
  \hline
  scid  & integer & subcorpus & starts from 1\\
  subcorpus & TEXT & short name & same across languages\\
  subtitle &TEXT & display name & \\
  url & TEXT & source & 
\\ \hline
\end{tabular}

  \begin{tabular}{llll}
  \textbf{scid}   &  \textbf{subcorpus} & \textbf{subtitle} & \textbf{url} \\
  \hline
  1 & spec & The Adventure of the Speckled Band &  www.??? \\
  2 & danc & The Adventure of the Dancing Men &  www.??? 
\end{tabular}

% e.g:  1,'sb', 'The Adventure of the Speckled Band', www.??? 

FIXME: maybe add license here?


\myslide{Monolingual Sentence Information}

We keep the original sentence.  \texttt{cfrom} and \texttt{cto} are calculated relative to this.

\texttt{pid} is currently unused.  We may need more META info (header, \ldots).


\noindent  \begin{tabular}{llll}
  \textbf{row}   &  \textbf{type} & \textbf{data} & \textbf{comment} \\\hline
  sid   & integer & sentence ID & unique over all corpora\\
  pid   & integer & paragraph ID & unique in subcorpus, start from 0\\
  sent  & text    & sentence & pre-tokenization\\
  scid  & integer & subcorpus & starts from 1\\
  comment & string & comment \\
\multicolumn{4}{l}{FOREIGN KEY(scid) REFERENCES subcorp(scid)}

\end{tabular}


\myslide{Monolingual Word Information}


\noindent \begin{tabular}{llll}
  \textbf{row}   &  \textbf{type} & \textbf{data} & \textbf{comment} \\\hline
  sid   & integer & sentence ID & unique over all corpora\\
  wid   & integer & word ID & unique in sentence, start from 0\\
  word  & string  & surface form &  \\
  pos   & string  & part of speech & per language  \\
  lemma & string  & lemma \\
  cfrom$^*$ & integer & cfrom & start character offset of word in sentence \\
  cto$^*$ & integer & cto & end character offset of word in sentence \\
  comment & string & comment \\
\hline
\end{tabular}
\\ $^*$ currently unused


\myslide{Monolingual Concept Information}

Currently cid is unique in a given corpus (genre).  When we link the
corpora together, this causes problems.

\noindent  \begin{tabular}{llll}
  \textbf{row}   &  \textbf{type} & \textbf{data} & \textbf{comment} \\\hline
  cid   & integer & concept ID  & unique in corpus, start from 1 \\ 
  sid   & integer & sentence ID & (sid, wid) gives the word \\
  wid   & integer & word ID & \ldots can have more than one\\
  clemma & string  & concept lemma & may be multi-word \\
  nlemma & string  & concept lemma & used for error in corpus \\
  tag   & string & wordnet id, tag, $\ldots$ \\
  ntag   & string & wordnet id, tag, $\ldots$ & used for tag not in tags\\
  tags & string & list of tags & for updating \\
  comment & string & comment \\
\hline
\end{tabular}


\myslide{Cross-Lingual Meta}
\begin{table}[htbp]
\centering
  \begin{tabular}{llll}
  \textbf{row}   &  \textbf{type} & \textbf{data} & \textbf{comment} \\
  \hline
  fcorpus & string & source corpus & name without \texttt{.db}\\
  tcorpus & string & target corpus & name without \texttt{.db}\\
  comment & string & comment \\
\hline
\end{tabular}
\caption{\texttt{meta}: Meta for links}
\end{table}

\myslide{SLINK Cross-Lingual Sentence Links}

\noindent  \begin{tabular}{llll}
  \textbf{row}   &  \textbf{type} & \textbf{data} & \textbf{comment} \\\hline
  slid & integer & link ID & unique in corpus pair \\
  fsid   & integer & source sentence ID & \\
  tsid   & integer & target sentence ID & \\
  ltype  & string  & link type &  \verb|=, ~, <, >|, Null if unknown\\
  conf   & float & confidence \\
  comment & string & comment \\
\hline
\end{tabular}


\myslide{WLINK Cross-Lingual Word Links}

Need to add \texttt{fsid}, \texttt{tsid} as wid numbers are unique only per sentence.

\noindent  \begin{tabular}{llll}
  \textbf{row}   &  \textbf{type} & \textbf{data} & \textbf{comment} \\\hline
  wlid & integer & link ID & unique in corpus pair \\
  fwid   & integer & source word ID & \\
  twid   & integer & target word ID & \\
  ltype  & string  & link type &  \verb|=, ~, <, >|\\
  conf   & float & confidence \\
  comment & string & comment \\
\hline
\end{tabular}

\myslide{CLINK Cross-Lingual Concept Information}

Need to add \fd{fsid}, \fd{tsid} when cid numbers become unique only per sentence.


\noindent  \begin{tabular}{llll}
  \textbf{row}   &  \textbf{type} & \textbf{data} & \textbf{comment} \\\hline
  clid & integer & link ID & unique in corpus pair \\
  fcid   & integer & source concept ID & \\
  tcid   & integer & target concept ID & \\
  ltype  & string  & link type &  \verb|=, ~, <, >|\\
  conf   & float & confidence \\
  comment & string & comment \\
\hline
\end{tabular}




\myslide{ToDo}

\begin{itemize}
\item Proper Back up
\item Small fixes (FCB and TTL)
\item Indonesian Quality Check: with lexical tagging (DM)
\item Chinese Quality Check: with lexical tagging (SW)
\item Japanese Quality Check: with lexical tagging (??)
\item English Quality Check: with lexical tagging (??)
\item Annotation Package
\item Wordnet tools
\item Decide on whether to tag unlexicalized MWEs
\item Add unknown words to wordnet and tag the remainder
\\ or allow partial tagging
\item Can also hire people
\end{itemize}




\end{CJK}
\clearpage
\end{document}


